\clearpage
\begin{examplebox}
\textbf{例 1｜最基础：先走两步，消掉两个负号}
\begin{center}\begin{tikzpicture}[font=\small,>=stealth,line join=round]
\draw[gray!40,rounded corners=3pt] (0,0) rectangle (14.2,8.3);\node at (7.1,7.700000000000001) {例1：每走2，高度在3与负3之间交替};
\begin{scope}[shift={(1.5444444444444445,3.828947368421052)},x=0.5555555555555555cm,y=0.4842105263157894cm]\draw[->,gray] (-0.8,0)--(10,0) node[right]{$x$};\draw[->,gray] (0,-4.5)--(0,5) node[above]{$y$};\draw[blue,thick,solid,domain=-0.4:9.5,samples=180] plot(\x,{3*sin(deg(pi*\x/2))});\fill[orange!85!black] (1,3) circle(2.2pt);\fill[orange!85!black] (3,-3) circle(2.2pt);\fill[orange!85!black] (5,3) circle(2.2pt);\fill[orange!85!black] (7,-3) circle(2.2pt);\fill[orange!85!black] (9,3) circle(2.2pt);\foreach \v in {1,3,5,7,9}{\node[below=5pt] at (\v,0){$\v$};}\end{scope}\node[anchor=west,align=left,text width=5.0cm] at (8.4,4.0) {$f=3\sin\frac{\pi x}2$\\[9pt]$1,3,5,7,9$\\[9pt]对应 $3,-3,3,-3,3$\\[9pt]隔4回原值\\[9pt]$f(9)=f(1)=3$};
\node[text width=12.6cm,align=center] at (7.1,.65) {示例满足初值3；不能把每次相反误认为每次相等。};
\end{tikzpicture}\end{center}



\textbf{题目：}已知 $f(x+2)=-f(x)$ 对每个实数成立，且 $f(1)=3$，求 $f(3)$、$f(5)$、$f(9)$。

\textbf{解：}每前进 $2$，函数值变号。因此
\[
\underbrace{f(1)}_{3},\quad
\underbrace{f(3)}_{-3},\quad
\underbrace{f(5)}_{3},\quad
\underbrace{f(7)}_{-3},\quad
\underbrace{f(9)}_{3}.
\]
所以 $\boxed{f(3)=-3,\quad f(5)=3,\quad f(9)=3}$。

\textbf{更快的算法：}周期可取 $4$；$9=1+2\times4$，所以 $f(9)=f(1)=3$。\textbf{看见差是偶数个 $a$，直接回原值；差是奇数个 $a$，直接变号。}

\medskip
\end{examplebox}\clearpage\begin{examplebox}\textbf{例 2｜一张区间表：先把大数折回去，再检查变号}
\begin{center}\begin{tikzpicture}[font=\small,>=stealth,line join=round]
\draw[gray!40,rounded corners=3pt] (0,0) rectangle (14.2,8.3);\node at (7.1,7.700000000000001) {例2：每段长3，下一段反号；接头不能同时取两值};
\begin{scope}[shift={(1.3891566265060242,3.828947368421052)},x=0.36144578313253006cm,y=0.4842105263157894cm]\draw[->,gray] (-0.8,0)--(15.8,0) node[right]{$x$};\draw[->,gray] (0,-4.5)--(0,5) node[above]{$y$};\draw[blue,thick,solid,domain=0:3,samples=180] plot(\x,{1*(\x-0)});\draw[blue,thick,solid,domain=3:6,samples=180] plot(\x,{-1*(\x-3)});\draw[blue,thick,solid,domain=6:9,samples=180] plot(\x,{1*(\x-6)});\draw[blue,thick,solid,domain=9:12,samples=180] plot(\x,{-1*(\x-9)});\draw[blue,thick,solid,domain=12:15,samples=180] plot(\x,{1*(\x-12)});\fill[orange!85!black] (0,0) circle(2.2pt);\fill[orange!85!black] (3,0) circle(2.2pt);\fill[orange!85!black] (6,0) circle(2.2pt);\fill[orange!85!black] (9,0) circle(2.2pt);\fill[orange!85!black] (12,0) circle(2.2pt);\fill[orange!85!black] (1,1) circle(2.2pt);\fill[orange!85!black] (10,-1) circle(2.2pt);\fill[orange!85!black] (13,1) circle(2.2pt);\draw[blue,fill=white] (3,3) circle(2.8pt);\draw[blue,fill=white] (6,-3) circle(2.8pt);\draw[blue,fill=white] (9,3) circle(2.8pt);\draw[blue,fill=white] (12,-3) circle(2.8pt);\draw[blue,fill=white] (15,3) circle(2.8pt);\fill[orange!85!black] (15,0) circle(2.2pt);\end{scope}\node[anchor=west,align=left,text width=5.0cm] at (8.4,4.0) {$[0,3)$ 上 $f=x$\\[9pt]$f(x+3)=-f(x)$\\[9pt]$f(1)=1,\ f(10)=-1$\\[9pt]$f(13)=1$\\[9pt]实点取到，空圈不取};
\node[text width=12.6cm,align=center] at (7.1,.65) {每段左端取0，右端的正负3不取；否则与变号条件发生冲突。};
\end{tikzpicture}\end{center}



\textbf{题目：}设 $f(x+3)=-f(x)$ 对任意实数成立，并且当 $0\le x<3$ 时 $f(x)=x$。求 $f(10)$、$f(13)$。

\textbf{解：}每增加 $3$ 变号；每增加 $6$ 还原。
\[
10=1+3\times3,\qquad 13=1+4\times3.
\]
先由已知区间得到 $f(1)=1$，故
\[
\begin{aligned}
f(10)&=f(1+3\cdot3)=-f(1)=\boxed{-1},\\
f(13)&=f(1+4\cdot3)=f(1)=\boxed{1}.
\end{aligned}
\]
\textbf{为什么区间特意写成 $[0,3)$？} 如果同时写“$0\le x\le3$ 时 $f(x)=x$”，就会得到 $f(0)=0$、$f(3)=3$；但题目又要求 $f(3)=-f(0)=0$，前后矛盾。\textbf{端点不是小事，题目条件必须能够同时成立。}

\medskip
\end{examplebox}\clearpage\begin{examplebox}\textbf{例 3｜高考常见：上千项直接配对相消}
\begin{center}\begin{tikzpicture}[font=\small,>=stealth,line join=round]
\draw[gray!40,rounded corners=3pt] (0,0) rectangle (14.2,8.3);\node at (7.1,7.700000000000001) {例3：相邻两项一上一下，函数值相加为0};
\begin{scope}[shift={(1.5918032786885248,3.7966666666666664)},x=0.9836065573770493cm,y=0.30666666666666664cm]\draw[->,gray] (-0.5,0)--(5.6,0) node[right]{$x$};\draw[->,gray] (0,-7)--(0,8) node[above]{$y$};\draw[blue,thick,solid,domain=-0.2:5.3,samples=180] plot(\x,{5*cos(deg(pi*\x))});\fill[orange!85!black] (0,5) circle(2.2pt);\fill[orange!85!black] (1,-5) circle(2.2pt);\fill[orange!85!black] (2,5) circle(2.2pt);\fill[orange!85!black] (3,-5) circle(2.2pt);\fill[orange!85!black] (4,5) circle(2.2pt);\fill[orange!85!black] (5,-5) circle(2.2pt);\foreach \v in {0,1,2,3,4,5}{\ifnum\v=0\node[below left=5pt] at (\v,0){$\v$};\else\node[below=5pt] at (\v,0){$\v$};\fi}\end{scope}\node[anchor=west,align=left,text width=5.0cm] at (8.4,4.0) {$f=5\cos\pi x,\ a=1$\\[9pt]偶数位置：5\\[9pt]奇数位置：$-5$\\[9pt]每两项相消\\[9pt]2026项为1013对};
\node[text width=12.6cm,align=center] at (7.1,.65) {图示前3对；任意满足条件的函数同样可以两项一组相消。};
\end{tikzpicture}\end{center}



\textbf{题目：}已知 $f(x+1)=-f(x)$，且 $f(0)=5$，求
\[
S=f(0)+f(1)+f(2)+\cdots+f(2025).
\]
\textbf{解：}把相邻两项看成一对：
\[
\underbrace{f(0)+f(1)}_{0}+
\underbrace{f(2)+f(3)}_{0}+\cdots+
\underbrace{f(2024)+f(2025)}_{0}.
\]
因为每对都等于 $0$，而 $0$ 到 $2025$ 一共 $2026$ 项，刚好 $1013$ 对，所以 $\boxed{S=0}$。

\textbf{眼前一亮：}本题其实\textbf{用不到} $f(0)=5$，不需要算出任何一项具体是多少；只要知道相邻两项相反就够了。

\medskip
\end{examplebox}\clearpage\begin{examplebox}\textbf{例 4｜辨别最小周期：等式说得少，不能自己加结论}
\begin{center}\begin{tikzpicture}[font=\small,>=stealth,line join=round]
\draw[gray!40,rounded corners=3pt] (0,0) rectangle (14.2,8.3);\node at (7.1,7.700000000000001) {例4：变号步长是π，重复却早在三分之二π出现};
\begin{scope}[shift={(1.4287671232876713,3.8842857142857143)},x=0.821917808219178cm,y=1.3142857142857143cm]\draw[->,gray] (-0.4,0)--(6.9,0) node[right]{$x$};\draw[->,gray] (0,-1.7)--(0,1.8) node[above]{$y$};\draw[blue,thick,solid,domain=-0.1:6.5,samples=180] plot(\x,{sin(deg(3*\x))});\fill[orange!85!black] (0.5235987755982988,1) circle(2.2pt);\fill[orange!85!black] (2.6179938779914944,1) circle(2.2pt);\fill[orange!85!black] (4.71238898038469,1) circle(2.2pt);\foreach \v/\lab in {2.094395/\frac{2\pi}3,4.18879/\frac{4\pi}3,6.283185/2\pi}{\node[below=5pt] at (\v,-1.25){$\lab$};}\end{scope}\node[anchor=west,align=left,text width=5.0cm] at (8.4,4.0) {$f=\sin3x,\ a=\pi$\\[9pt]$f(x+\pi)=-f(x)$\\[9pt]$2\pi$ 是一个周期\\[9pt]最小正周期 $\frac{2\pi}{3}$};
\node[text width=12.6cm,align=center] at (7.1,.65) {三个波形在长度2π内重复；两次变号只保证得到一个周期。};
\end{tikzpicture}\end{center}



\textbf{题目：}若 $f(x+\pi)=-f(x)$，能否断定最小正周期是 $2\pi$？

\textbf{解：}不能。取 $f(x)=\sin3x$，则
\[
f(x+\pi)=\sin(3x+3\pi)=-\sin3x=-f(x),
\]
但 $\sin3x$ 的最小正周期是 $\frac{2\pi}{3}$，比 $2\pi$ 小。

\textbf{结论：}$2\pi$ 是它\textbf{一个}正周期，却不一定是\textbf{最小}正周期。

\medskip
\end{examplebox}\clearpage\begin{examplebox}\textbf{例 5｜附加奇偶性后才谈对称}
\begin{center}\begin{tikzpicture}[font=\small,>=stealth,line join=round]
\draw[gray!40,rounded corners=3pt] (0,0) rectangle (14.2,8.5);\node at (7.1,7.9) {同样加π变号，额外的奇偶性决定不同对称};
\node at (2.4,6.9) {$f=\sin x$：奇函数};\begin{scope}[shift={(1.2666666666666666,4.341935483870968)},x=0.7083333333333334cm,y=1.1612903225806452cm]\draw[->,gray] (-0.8,0)--(4,0) node[right]{$x$};\draw[->,gray] (0,-1.5)--(0,1.6) node[above]{$y$};\draw[blue,thick,solid,domain=-0.5:3.8,samples=160] plot(\x,{sin(deg(\x))});\fill[orange!85!black] (0.5235987755982988,0.5) circle(2pt);\fill[orange!85!black] (2.6179938779914944,0.5) circle(2pt);\draw[orange!85!black,dashed] (1.570796,-1.2)--(1.570796,1.3);\end{scope}\node[align=center,text width=4cm] at (2.4,1.7) {关于直线 $x=\frac\pi2$ 对称};\node at (7.050000000000001,6.9) {$g=\cos x$：偶函数};\begin{scope}[shift={(5.916666666666667,4.341935483870968)},x=0.7083333333333334cm,y=1.1612903225806452cm]\draw[->,gray] (-0.8,0)--(4,0) node[right]{$x$};\draw[->,gray] (0,-1.5)--(0,1.6) node[above]{$y$};\draw[violet,thick,solid,domain=-0.5:3.8,samples=160] plot(\x,{cos(deg(\x))});\fill[orange!85!black] (0.5235987755982988,0.8660254037844386) circle(2pt);\fill[orange!85!black] (2.6179938779914944,-0.8660254037844386) circle(2pt);\fill[orange!85!black] (1.570796,0) circle(2.2pt);\draw[orange!85!black,dashed] (.523599,.866025)--(2.617994,-.866025);\end{scope}\node[align=center,text width=4cm] at (7.050000000000001,1.7) {关于点 $(\frac\pi2,0)$\\中心对称};\node at (11.7,6.9) {$f=\sin x$ 的中心检验};\begin{scope}[shift={(10.566666666666666,4.341935483870968)},x=0.7083333333333334cm,y=1.1612903225806452cm]\draw[->,gray] (-0.8,0)--(4,0) node[right]{$x$};\draw[->,gray] (0,-1.5)--(0,1.6) node[above]{$y$};\draw[blue,thick,solid,domain=-0.5:3.8,samples=160] plot(\x,{sin(deg(\x))});\fill[orange!85!black] (1.5707963267948966,1) circle(2pt);\end{scope}\node[align=center,text width=4cm] at (11.7,1.7) {$f(\frac\pi2)=1\ne0$\\$(\frac\pi2,0)$ 不是中心};
\node[text width=12.6cm,align=center] at (7.1,.65) {三图同尺度；只有变号条件时，不能直接指定对称轴或对称中心。};
\end{tikzpicture}\end{center}



\textbf{题目：}已知 $f(x+\pi)=-f(x)$。分别讨论还知道“$f$ 为奇函数”或“$f$ 为偶函数”时能得到什么。

\textbf{解：}
\begin{itemize}[leftmargin=*,itemsep=0.35em]
\item 若 $f$ 是\textbf{奇函数}，则 $f(\frac\pi2+t)=f(\frac\pi2-t)$，图像关于直线 $x=\frac\pi2$ 对称。$\sin x$ 就是例子。
\item 若 $f$ 是\textbf{偶函数}，则 $f(\frac\pi2+t)=-f(\frac\pi2-t)$，图像关于点 $(\frac\pi2,0)$ 中心对称。$\cos x$ 就是例子。
\end{itemize}
\textbf{注意：}这两条都需要额外的奇偶性条件；只有“加 $\pi$ 变号”时不能任意选一条。
\end{examplebox}
\clearpage
\begin{examplebox}
\textbf{例 6｜最快识别错位反号：先相减，再乘二}
\begin{center}\begin{tikzpicture}[font=\small,>=stealth,line join=round]
\draw[gray!40,rounded corners=3pt] (0,0) rectangle (14.2,7.4);
\node[text width=12.6cm,align=center] at (7.1,6.75) {为什么要相减：看输入在数轴上的距离};

\draw[->,gray] (.8,3.5)--(13.1,3.5) node[right]{$u$};
\foreach \u/\lab in {1.4/x,5.4/{x+2},11.4/{x+5}}{
\draw (\u,3.4)--(\u,3.6);\node[below=6pt] at (\u,3.4){$\lab$};}
\draw[<->,teal,thick] (1.4,4.2)--(5.4,4.2) node[midway,above]{$2$};
\draw[<->,blue,thick] (5.4,4.2)--(11.4,4.2) node[midway,above]{$5-2=3$};
\draw[<->,violet] (1.4,5.3)--(11.4,5.3) node[midway,above]{$5$};
\node[text width=11.8cm,align=center] at (7.1,1.9) {输入相距 $3$：先变号；再走 $3$：恢复。\\$5+2=7$ 把两段从 $x$ 出发的距离相加，不能表示这两点的间距。};

\node[text width=12.6cm,align=center] at (7.1,.65) {数轴按比例绘制；$f(x+5)=-f(x+2)$ 给出的一个正周期为 $6$。};
\end{tikzpicture}\end{center}


\textbf{题目：}函数定义在 $\mathbb R$ 上，对每个实数 $x$ 有 $f(x+5)=-f(x+2)$。求一个正周期。

\textbf{解：}令 $t=x+2$，则 $x+5=t+3$，所以
\[f(t+3)=-f(t),\qquad f(t+6)=-f(t+3)=f(t).\]
故 $\boxed{6\text{ 是一个正周期}}$。最快识别为 $2|5-2|=6$，不能仅凭题设断定最小。

\textbf{回验条件可以成立：}取 $f(u)=\cos(\pi u/3)$，则 $f(x+5)=\cos(\pi(x+2)/3+\pi)=-f(x+2)$，再移 $6$ 正好重复。这个例子只作回验，前面的代数证明才适用于所有满足条件的函数。

\end{examplebox}
